\subsection{Critical checks for the miscatalogued exterior-algebra lecture}

These checks concern record~437, whose manuscript date is 13 February 1978.
They audit the supplied material without turning it into evidence of a
1976--77 research problem. The results are elementary verifications of the
construction recorded in the notebook; no claim of novelty is intended.

\begin{proposition}[The regular exterior module is acyclic]\ours{}
Let \(k\) be a field, let \(\Xi\) have finite dimension \(m\geq1\), and
let \(0\ne\xi\in\Xi\). Give \(\Lambda=\bigwedge\Xi\) its negative
grading and differential \(d(v)=\xi\wedge v\). Then the complex is
contractible. Consequently every finite direct sum of grading shifts of
\(\Lambda\) gives the zero bundle under the construction of record~437.
\end{proposition}
\begin{proof}
Choose \(a\in\Xi^*\) with \(a(\xi)=1\). Define contraction on a simple
exterior tensor by
\[
 \iota_a(v_1\wedge\cdots\wedge v_r)
 =\sum_{j=1}^r(-1)^{j-1}a(v_j)
 v_1\wedge\cdots\wedge\widehat{v_j}\wedge\cdots\wedge v_r.
\]
Applying this formula to \(\xi\wedge v\), the term contracting the first
factor is \(v\); each remaining term is the negative of the corresponding
term of \(\xi\wedge\iota_a(v)\). Therefore
\(\iota_a d+d\iota_a=1\). If \(dv=0\), this identity gives
\(v=d\iota_a(v)\), so every cycle is a boundary. The homotopy commutes
with taking direct sums and does not depend on the choice of grading
origin. In particular the output has rank zero, which resolves the
notebook's ambiguous use of ``trivial'' in this example.
\end{proof}

\begin{remark}[Independent checks]\oursComputed{}
The identity was independently re-derived by writing
\(\Xi=k\xi\oplus U\), hence
\(\Lambda=\bigwedge U\oplus\xi\wedge\bigwedge U\). The differential
maps the first summand isomorphically to the second and kills the second.
The exact-arithmetic program \texttt{code/critical\_checks.py} additionally
constructs exterior multiplication and contraction by two separate
bit-mask formulas. For \(m=1,\ldots,6\) and three nonzero vectors in each
dimension, all 18 cases satisfy \(d^2=0\), \(\iota_a d+d\iota_a=1\),
and zero homology in every degree. The ranks agree with
\(\binom{m-1}{r}\) in exterior degree \(r\). These finite checks support
the implementation; the preceding proof establishes the assertion for
every \(m\), not an extrapolation from the computed cases.
\end{remark}

\begin{example}[The inverse construction requires a grading convention]\ours{}
Take \(n=0\), so the base is a point, and let
\(\mathcal E=\mathcal O\). For every \(N\geq0\), the two consecutive
section spaces are one-dimensional and \(X_0\) is the identity. If the
two tensor factors in P000756 are given their raw exterior grading, then
\[
 \partial:\Lambda\longrightarrow\Lambda,\qquad
 \partial(1)=\xi_0,\qquad \partial(\xi_0)=0.
\]
Its kernel is \(k\xi_0\), concentrated in internal degree \(-1\).
The action of \(\xi_0\) on that kernel is zero, so its fibre homology
is nonzero in degree \(-1\); the raw kernel is not proper. This is an
explicit check of an omitted grading convention, not a counterexample to
the BGG correspondence. The dual free module
\(\operatorname{Hom}_k(\Lambda,k)\) differs from \(\Lambda\) by a
grading shift, which is essential in the precise inverse construction
\citep[\S4]{bernsteingelfandgelfand1978bundles}.
\end{example}

\begin{proposition}[Sign of the twist operation]\ours{}
Use the negative grading of record~437. If \(V\) is proper and
\(0\to V\to P\to V'\to0\) is an exact sequence of finite graded
exterior-algebra modules with \(P\) free, set \(W_j=V'_{j+1}\).
Then \(W\) is proper and \(\Phi(W)=\Phi(V)(1)\).
\end{proposition}
\begin{proof}
For each nonzero \(\xi\), applying \(L_\xi\) gives a short exact
sequence of complexes. The middle complex is acyclic by the proposition
above, so its homology sequence gives
\(H_j(L_\xi(V'))\cong H_{j-1}(L_\xi(V))\). Consequently the only
possibly nonzero homology of \(V'\) is in chain degree one, and the
displayed grading shift moves it to degree zero. On projective space,
\[
 \mathcal L^j(W)
 =V'_{1-j}\otimes\mathcal O(j)
 =\mathcal L^{j-1}(V')(1).
\]
The short exact sequence of sheaf complexes gives
\(\mathcal H^{-1}(\mathcal L(V'))\cong\mathcal H^0(\mathcal L(V))\).
Reindexing by the displayed formula proves the claimed twist. Depending
on the convention for a shifted complex, a minus sign in its differential
is removed by multiplying degree \(j\) by \((-1)^j\); cohomology and
the bundle are unaffected.
\end{proof}

\begin{remark}[Second check of indexing]
There are two independent index tests: the chain homology sequence
places the cokernel's fibre homology in degree one, whereas the sheaf
cohomology sequence places it in degree minus one. The substitution
\(j\mapsto-j\) in \(\mathcal L^j(V)=V_{-j}\otimes\mathcal O(j)\)
reconciles them and gives \(\mathcal O(1)\), not \(\mathcal O(-1)\).
This is exactly the convention in remark~2 after Theorem~1 of
\citet{bernsteingelfandgelfand1978bundles}.
\end{remark}


\begin{remark}[Independent check of the inverse-module indexing]\ours{}
The conventions of occurrence~1977-437.19 were checked against
\citet[\S2, Corollary~2.4, Theorem~3.7 and \S4]{eisenbudfloystadschreyer2003exterior}.
Put \(Z_m=\ker(T^m\to T^{m+1})\), use internal shifts
\(Q(a)_j=Q_{a+j}\), and use cochain shifts \(K[a]^j=K^{a+j}\).
At a sufficiently large cutoff \(N\), the BGG reciprocity theorem,
applied to the regraded kernel, gives
\(\mathcal L(Z_N(N))\simeq\mathcal E(N)\).
Direct substitution in the definition of the linear complex gives
\[
 \mathcal L(Q(N))=\mathcal L(Q)[-N](N).
\]
Therefore \(\mathcal L(Z_N)\simeq\mathcal E[N]\).
For every integer \(m\), exactness of the Tate complex gives
\(0\to Z_m\to T^m\to Z_{m+1}\to0\).
The middle term is graded free, and its linear sheaf complex is acyclic
by the contraction above. Its associated triangle therefore gives
\(\mathcal L(Z_{m+1})\simeq\mathcal L(Z_m)[1]\).
Iteration yields \(\mathcal L(Z_N)\simeq\mathcal L(Z_0)[N]\), hence
\(\mathcal L(Z_0)\simeq\mathcal E\) with no additional twist.
The bounded complex of locally free sheaves may be split locally away
from its locally free degree-zero cohomology, so its fibres have no
other homology: \(Z_0\) is proper. The point case gives the trivial
exterior module in internal degree zero. On \(\mathbb P^1\), the case
\(\mathcal E=\mathcal O(-1)\) gives the Euler kernel of
\(\mathcal O^2\to\mathcal O(1)\), again checking the grading sign.
These are checks of conventions in the cited construction, not a new
proof of the BGG reciprocity theorem.
\end{remark}
